\documentclass[12pt]{article}
\usepackage{amsmath,amssymb,geometry,enumitem}
\geometry{margin=1in}
\setlength{\parskip}{0.6em}
\setlength{\parindent}{0pt}

\begin{document}

\begin{center}
{\LARGE \textbf{ECON 300 -- Labour Economics}}\\[0.4em]
{\Large \textbf{Solutions: Final Exam Numerical Problems }}\\[0.4em]

\end{center}

\vspace{0.5em}

% ============================================================
\section*{Question 1: Labour Supply, Policy, and Unemployment}
% ============================================================

Johny has:
\[
T=100,\qquad h=T-L,\qquad w=20,\qquad V=200
\]
where \(L\) is leisure, \(h\) is hours worked, and \(V\) is non-labour income.

\begin{enumerate}[label=(\alph*)]

\item \textbf{Budget constraint}

Income is labour earnings plus non-labour income:
\[
C=wh+V
\]

Since:
\[
h=T-L=100-L
\]

Substitute into the income equation:
\[
C=20(100-L)+200
\]

\[
C=2000-20L+200
\]

\[
\boxed{C=2200-20L}
\]

\textbf{Income intercept:}

Set \(L=0\):
\[
C=2200-20(0)=2200
\]

\[
\boxed{2200}
\]

\textbf{Leisure intercept:}

Set hours worked equal to zero, so \(L=T=100\).

\[
\boxed{100}
\]

\textbf{Slope:}

From
\[
C=2200-20L
\]
the slope is:

\[
\boxed{-20}
\]

Interpretation: one more hour of leisure reduces income by \(\$20\).

% ------------------------------------------------------------

\item \textbf{Guaranteed income program}

Given:
\[
G=600,\qquad t=0.5
\]

Disposable income while benefits are positive is:
\[
C=G+(1-t)wh
\]

Substitute:
\[
C=600+(1-0.5)(20)h
\]

\[
C=600+0.5(20)h
\]

\[
\boxed{C=600+10h}
\]

\textbf{Effective net wage:}
\[
(1-t)w=0.5(20)=10
\]

\[
\boxed{\$10}
\]

\textbf{Break-even level of weekly earnings:}

Benefits go to zero when:
\[
G=twh
\]
or:
\[
wh=\frac{G}{t}
\]

\[
wh=\frac{600}{0.5}=1200
\]

\[
\boxed{\$1200}
\]

\textbf{Break-even hours worked:}

Since \(w=20\):
\[
20h=1200
\]

\[
h=60
\]

\[
\boxed{60\text{ hours}}
\]

% ------------------------------------------------------------

\item \textbf{Change in labour supply}

Initially Johny works:
\[
40\text{ hours}
\]

After the program:
\[
30\text{ hours}
\]

Change in labour supply:
\[
30-40=-10
\]

\[
\boxed{-10\text{ hours}}
\]

So labour supply falls by 10 hours.

\textbf{Interpretation:}

This is consistent with:
\begin{itemize}
    \item an \textbf{income effect}: the guaranteed income makes Johny richer even without working, so he may choose more leisure;
    \item a \textbf{substitution effect}: the effective wage falls from \(\$20\) to \(\$10\), so leisure becomes cheaper relative to work.
\end{itemize}

Both effects reduce labour supply.

% ------------------------------------------------------------

\item \textbf{Labour force, unemployed, employed}

Working-age population:
\[
5{,}000{,}000
\]

LFPR:
\[
70\%
\]

Unemployment rate:
\[
6\%
\]

\textbf{Labour force:}
\[
0.70(5{,}000{,}000)=3{,}500{,}000
\]

\[
\boxed{3{,}500{,}000}
\]

\textbf{Unemployed workers:}
\[
0.06(3{,}500{,}000)=210{,}000
\]

\[
\boxed{210{,}000}
\]

\textbf{Employed workers:}
\[
3{,}500{,}000-210{,}000=3{,}290{,}000
\]

\[
\boxed{3{,}290{,}000}
\]

% ------------------------------------------------------------

\item \textbf{New unemployment and labour force}

Initial unemployed:
\[
210{,}000
\]

Initial labour force:
\[
3{,}500{,}000
\]

Changes:
\begin{itemize}
    \item 60,000 unemployed find jobs \(\Rightarrow\) unemployment falls
    \item 40,000 employed lose jobs and search \(\Rightarrow\) unemployment rises
    \item 30,000 workers leave the labour force
\end{itemize}

\textbf{New unemployed:}

Start with initial unemployed:
\[
210{,}000
\]

Subtract those who find jobs:
\[
210{,}000-60{,}000=150{,}000
\]

Add those who lose jobs:
\[
150{,}000+40{,}000=190{,}000
\]

\[
\boxed{190{,}000}
\]

\textbf{New labour force:}

The 30,000 workers leaving the labour force reduce the labour force:

\[
3{,}500{,}000-30{,}000=3{,}470{,}000
\]

\[
\boxed{3{,}470{,}000}
\]

\textbf{New unemployment rate:}

\[
\frac{190{,}000}{3{,}470{,}000}\times 100 \approx 5.48\%
\]

\[
\boxed{5.48\%}
\]

\textbf{Interpretation:}

Even though some workers lost jobs, the unemployment rate falls because many unemployed found jobs and some workers left the labour force.

\end{enumerate}

% ============================================================
\section*{Question 2: Labour Demand, Monopsony, Immigration, and Human Capital}
% ============================================================

Demand:
\[
W=50-0.5N
\]

Supply:
\[
W=10+0.5N
\]

\begin{enumerate}[label=(\alph*)]

\item \textbf{Competitive equilibrium}

Set demand equal to supply:
\[
50-0.5N=10+0.5N
\]

\[
40=N
\]

\[
\boxed{N=40}
\]

Substitute into either curve:
\[
W=10+0.5(40)=10+20=30
\]

\[
\boxed{W=30}
\]

So the competitive equilibrium is:
\[
\boxed{(N,W)=(40,30)}
\]

% ------------------------------------------------------------

\item \textbf{Monopsony hiring 20 workers}

If the monopsonist hires:
\[
N=20
\]

Then the wage paid comes from the supply curve:
\[
W=10+0.5(20)=10+10=20
\]

\[
\boxed{W=20}
\]

Compare this to competitive wage:
\[
20<30
\]

\[
\boxed{\text{Monopsony wage is \$10 lower than the competitive wage}}
\]

\textbf{Explanation:}

A monopsonist has labour market power and faces an upward-sloping supply curve. It can lower wages by restricting employment below the competitive level.

% ------------------------------------------------------------

\item \textbf{Minimum wage of \$30}

\textbf{Workers willing to work at \$30:}

Use the supply curve:
\[
30=10+0.5N
\]

\[
20=0.5N
\]

\[
N=40
\]

\[
\boxed{40}
\]

\textbf{Workers demanded at \$30:}

Use the demand curve:
\[
30=50-0.5N
\]

\[
0.5N=20
\]

\[
N=40
\]

\[
\boxed{40}
\]

\textbf{Actual employment:}

Employment is the smaller of labour supply and labour demand:
\[
\boxed{40}
\]

So:
\[
\boxed{\text{Actual employment}=40}
\]

% ------------------------------------------------------------

\item \textbf{Immigration shares}

Total immigrants:
\[
400{,}000
\]

\textbf{Economic share:}
\[
\frac{200{,}000}{400{,}000}\times 100=50\%
\]

\[
\boxed{50\%}
\]

\textbf{Family share:}
\[
\frac{120{,}000}{400{,}000}\times 100=30\%
\]

\[
\boxed{30\%}
\]

\textbf{Refugee share:}
\[
\frac{80{,}000}{400{,}000}\times 100=20\%
\]

\[
\boxed{20\%}
\]

If economic immigrants rise by 10\%:
\[
200{,}000(1.10)=220{,}000
\]

New total:
\[
220{,}000+120{,}000+80{,}000=420{,}000
\]

\[
\boxed{420{,}000}
\]

% ------------------------------------------------------------

\item \textbf{Human capital investment}

Given:
\[
A=3000,\qquad r=0.05,\qquad T=8
\]

\textbf{Present value of benefits:}
\[
PV=3000\left(\frac{1-(1.05)^{-8}}{0.05}\right)
\]

First compute:
\[
(1.05)^{-8}\approx 0.67684
\]

So:
\[
PV=3000\left(\frac{1-0.67684}{0.05}\right)
\]

\[
PV=3000\left(\frac{0.32316}{0.05}\right)
\]

\[
PV=3000(6.4632)=19{,}389.60
\]

\[
\boxed{PV\approx \$19{,}389.60}
\]

\textbf{Net present value:}
\[
NPV=19{,}389.60-12{,}000=7{,}389.60
\]

\[
\boxed{NPV\approx \$7{,}389.60}
\]

\textbf{Decision:}

Because:
\[
NPV>0
\]
Johny should invest.

\[
\boxed{\text{Yes, Johny should invest}}
\]

\end{enumerate}

\end{document}